\section{Method}
\label{sec:method}

Given two source assets and an interpolation coefficient $\alpha\in[0,1]$,
\methodname{} generates a textured intermediate asset $\widehat M_\alpha$.
We represent the sources as $A=(P_A,F_A)$ and $B=(P_B,F_B)$ in TRELLIS
SLat space~\cite{xiang2025trellis}, where $P_k$ contains active voxel
coordinates and $F_k$ their latent features.
The inputs can be cached SLat representations or existing meshes encoded
into this space. Their geometry is therefore provided to the morphing
model, rather than reconstructed from endpoint images.
The convention $\alpha=1$ selects $A$ and $\alpha=0$ selects $B$.

\subsection{Encoding the Source Geometry}
\label{sec:representation}

To make this geometry available to the generator, we use a sparse 3d convolutional encoder.
An input projection is followed by residual convolutional blocks and two
stride-two convolutions, reducing the spatial resolution from $64^3$ to
$16^3$.
Each residual block contains two $3\times3\times3$ sparse convolutions with
group normalization and SiLU activations, and an additive skip connection
that preserves the block input at the same resolution.
A final convolutional projection produces the spatial descriptors.
Positional embeddings retain their location within the 3D grid before
projection and normalization to the flow model's conditioning width.

Within each generative stage $s$, the same encoder $E_s$ processes the two
sources separately, producing $H_A^s=E_s(P_A,F_A)$ and
$H_B^s=E_s(P_B,F_B)$.
Keeping the two conditions distinct lets the generator control their
contributions throughout the morph instead of mixing the source latents
before generation.

\subsection{Generating the Intermediate Asset}
\label{sec:source_fusion}

These conditions guide both stages of the TRELLIS pipeline.
The sparse-structure (SS) flow generates a coarse latent, which its decoder
converts to active voxel coordinates. The SLat flow then generates the
features attached to those coordinates.
We initialize both flow models with pretrained TRELLIS weights and adapt
them using LoRA~\cite{hu2022lora}; the new source encoders and fusion gates
are learned jointly with the adapters.
The original backbone weights and asset decoders remain frozen.

Writing $\mathbf c_s=(H_A^s,H_B^s,\alpha)$, generation proceeds as
\begin{equation}
\begin{aligned}
\widehat S_\alpha
  &=G_{\mathrm{SS}}(\epsilon_{\mathrm{SS}};\mathbf c_{\mathrm{SS}}),\\
\widehat P_\alpha
  &=D_{\mathrm{SS}}(\widehat S_\alpha),\\
\widehat F_\alpha
  &=G_{\mathrm{SL}}(\epsilon_{\mathrm{SL}};
                   \widehat P_\alpha,\mathbf c_{\mathrm{SL}}),\\
\widehat M_\alpha
  &=D_{\mathrm{SL}}(\widehat P_\alpha,\widehat F_\alpha).
\end{aligned}
\label{eq:method_pipeline}
\end{equation}
Here $G_s$ samples the stage's flow from Gaussian noise $\epsilon_s$,
$D_{\mathrm{SS}}$ decodes the active support, and $D_{\mathrm{SL}}$ decodes
the final asset. SLat sampling takes place on the support $\widehat P_\alpha$
predicted by the first stage.

To control how each source guides this generation, every transformer block
evaluates cross-attention separately for $H_A^s$ and $H_B^s$.
The resulting responses $h_A,h_B$ are combined by a learned gate:
\begin{equation}
\begin{aligned}
r_\phi&=\operatorname{MLP}_\phi(h_A,h_B,\alpha),\\
g&=(1-\alpha)+\alpha(1-\alpha)\tanh(r_\phi),\\
h_{\mathrm{mix}}&=(1-g)\odot h_A+g\odot h_B.
\end{aligned}
\label{eq:method_gate}
\end{equation}
The gate varies across tokens and channels, allowing different parts of
the representation to draw differently from the two sources.
A zero-initialized output layer starts from linear fusion.
The learned correction vanishes at both endpoints, so the selected source
alone contributes to the attention response there.

\subsection{Learning from the Teacher and the Prior}
\label{sec:synthetic_training}

To learn this generator without annotated transitions, we run
Interp3D~\cite{liu2026interp3d} on image pairs and capture
its intermediate SS latents and SLat features. These provide targets for the two stages, while the
student receives the SLat representations of the source assets.
Four interpolation coefficients per pair cover different parts of the
transition; a small fraction of target slots instead uses the exact source
latents to supervise endpoint reconstruction.
We also exchange the sources and replace $\alpha$ with $1-\alpha$, reusing
the same target to encourage source-order consistency.

We train the two stages separately with standard conditional flow
matching~\cite{lipman2022flowmatching}, denoting the corresponding
teacher-supervised loss by $\mathcal L_{\mathrm{FM}}$.
SLat training uses the target's voxel support, whereas inference uses the
support predicted by the SS stage.

Teacher supervision alone can transfer artifacts to the student.
We therefore also sample a latent from the student during training and
use the original, frozen TRELLIS model as a generative guide.
After a small noise perturbation, TRELLIS estimates corrections under
each endpoint image condition.
We limit these corrections and weight their contributions by $\alpha$,
providing corrected references without replacing the learned morph.
The prior loss encourages the student sample to approach these references
and discourages unstable latent statistics.
The corrected references are treated as fixed targets: gradients pass
through the final student sampling steps to update the LoRA adapters,
source encoders, and gates, not through TRELLIS or its target construction.
Thus, the prior changes how the student generates future samples rather
than simply editing its current output.

The resulting objective for each stage is
\begin{equation}
\mathcal L_{\mathrm{total}}
=\mathcal L_{\mathrm{FM}}
+\lambda_{\mathrm{prior}}\mathcal L_{\mathrm{prior}},
\label{eq:method_objective}
\end{equation}
where $\lambda_{\mathrm{prior}}$ balances teacher supervision and sample-level
regularization. 
Once the source SLat representations are available, neither the teacher,
the endpoint images, nor the auxiliary prior is needed for morphing inference.